\documentclass[11pt]{article}
\usepackage[a5paper,margin=16mm]{geometry}
\usepackage{fontspec}
\setmainfont{Latin Modern Roman}
\setsansfont{Latin Modern Sans}
\usepackage{xcolor}
\usepackage{tikz}
\usetikzlibrary{arrows.meta,tikzmark,calc,fit,shapes.misc}
\newcommand{\hl}[2]{\tikz[remember picture,baseline=(#1.base)]\node[inner sep=0pt,outer sep=0pt,anchor=base] (#1) {#2};}

\begin{document}
\section*{Annotating prose with remembered positions}
A \hl{subj}{compiler} translates a \hl{obj}{program} into another language, while an
\hl{int}{interpreter} executes it directly. Every \hl{tool}{toolchain} combines them,
so the arrows below can only be placed once the second pass knows where each word landed.

\begin{tikzpicture}[remember picture,overlay,>={Stealth[length=2mm]},every node/.style={font=\sffamily\scriptsize}]
  \draw[->,thick,blue] (subj.north) to[bend left=35] node[above] {translates} (obj.north);
  \draw[->,thick,red!70!black] (int.south) to[bend right=35] node[below] {executes} (obj.south);
  \node[draw,rounded corners,fill=yellow!20,fit=(tool),inner sep=1.5pt] {};
\end{tikzpicture}

\bigskip
Second paragraph with a \tikzmark{start}tikzmark start, filler text that runs for a
while to force the wrapped line to be longer than one full line of the narrow measure
here, and the \tikzmark{stop}tikzmark stop.

\begin{tikzpicture}[remember picture,overlay]
  \fill[green!25,opacity=.5] ($(pic cs:start)+(-1pt,-2pt)$) rectangle ($(pic cs:stop)+(1pt,9pt)$);
  \node[anchor=west,font=\sffamily\scriptsize,text=green!40!black] at ($(pic cs:stop)+(4mm,0)$) {span end};
\end{tikzpicture}

\vspace{2.2cm}
The \texttt{remember picture} option writes node positions to the \texttt{.aux}
file; this paragraph sits below two overlays that have been resolved.
\end{document}
